\subsection{A weight-tied stage as an autonomous control}
\label{sec:weight_tied}

To separate what depends on having nine independent blocks from what depends only on
the residual form, we train a ConvNeXt-T whose stage~3 applies one shared block nine
times.
Reusing blocks of a model that was trained with distinct blocks tests a different
object, and fails (\cref{app:block-ops}); here the stage is \emph{trained} to use a
single function.
Both models are trained by us on ImageNet-1k for $300$ epochs with the same
augmentation, regularisation, optimiser and effective batch size of $1024$.
The runs are not matched in every setting: the untied model uses peak learning rate
$5{\times}10^{-4}$ with a $20$-epoch warmup on a single process, the tied one
$1{\times}10^{-3}$ with a $10$-epoch warmup on four distributed processes; no run
fixes a seed (\cref{app:training}).
The accuracy gap below is therefore indicative, not a controlled ablation.

The untied ConvNeXt-T reaches $80.96\%$ top-1 with $28.59$M parameters, below the
$82.1\%$ of the original recipe~\citep{liu2022convnext}, which uses a batch of $4096$.
The weight-tied model reaches $80.34\%$ with $18.97$M parameters ($33\%$ fewer,
$-0.62$ points), and sharing one block per stage in all four stages reaches $78.99\%$
with $8.68$M ($70\%$ fewer, $-1.97$ points).
One residual function applied nine times is thus close to nine independent ones.
The tied stage is a clean autonomous ODE $\dot x=f(x)$ over $T{=}9$, and we use it as
the control against which the switched stack is read.
